\section{Технический проект}
\subsection{Словарь предметной области}

На основании анализа предметной области был составлен словарь предметной области. В словаре представлены термины на русском и английском языках.

В таблице \ref{dictionary:table} представлен словарь предметной области для программной части разрабатываемого программного продукта.

\renewcommand{\arraystretch}{0.8} % уменьшение расстояний до сетки таблицы
\begin{xltabular}{\textwidth}{|p{3.6cm}|X|}
	\caption{Словарь предметной области разрабатываемой программной системы\label{dictionary:table}}\\
	\hline \centrow \setlength{\baselineskip}{0.7\baselineskip} Термин & \centrow \setlength{\baselineskip}{0.7\baselineskip} Описание \\
	\hline \centrow 1 & \centrow 2\\ \hline
	\endfirsthead
	\caption*{Продолжение таблицы \ref{dictionary:table}}\\
	\hline \centrow 1 & \centrow 2\\ \hline
	\finishhead
	QMK (Quantum Mechanical Keyboard) & Открытое программное обеспечение (прошивка) для механической клавиатуры, позволяющее полностью контролировать и настраивать каждое действие, связанное с клавишами, слоями, подсветкой и другим функционалом. Основное преимущество QMK -- возможность глубокой кастомизации клавиатуры без необходимости знания программирования. \\ \hline
	Ключ-код (Keycode) & Специальный идентификатор, присваиваемый физической клавише на клавиатуре. Через ключ-коды определяются действия, происходящие при нажатии той или иной клавиши. Например, стандартный ключ-код для клавиши "<A"> -- это KC\_A. \\ \hline
	Веб-приложение (Web Application ) & Прикладное программное обеспечение, которое работает на веб-сервере, в отличие от компьютерных программ, которые запускаются локально в операционной системе устройства. \\ \hline
	Слой (Layer) & Система, позволяющая клавиатуре обрабатывать разные функции для одной и той же клавиши в зависимости от текущего активного слоя. Переключение между слоями осуществляется специальными клавишами или сочетаниями. Например, один слой может отвечать за стандартные символы, другой -- за мультимедийные команды, третий -- за макросы и т.п. \\ \hline
	Макрос (Macro) & Последовательность действий, задаваемых на определенную клавишу или комбинацию клавиш. Макросы могут эмулировать нажатие нескольких клавиш, запускать программы, вводить текст и многое другое. \\ \hline
	Модификатор (Modifier) & Тип клавиш, изменяющий поведение других клавиш при их совместном нажатии. Модификаторами являются такие клавиши, как SHIFT, CTRL, ALT и др., изменяющие результат нажатия обычных клавиш (например, сочетание SHIFT+A даст заглавную букву "<A">). \\ \hline
	Комбинация (Combo) & Возможность привязать одно действие к двум или более клавишам, нажатым одновременно. Часто используется для экономии физического пространства на клавиатуре, позволяя назначить редкие действия на сочетание редких клавиш. \\ \hline
	Динамическая смена слоев (Dynamic Layer Change) & Способность переключаться между слоями без постоянного удерживания специальной клавиши (обычная ситуация -- MO()). При включении динамической смены слои меняются по однократному нажатию определенных клавиш. \\ \hline
	Режим удержания (Tap-Hold Mode) & Режим, в котором реакция клавиши зависит от продолжительности её нажатия. Например, короткое нажатие вызывает одно действие, длительное удержание -- другое (это характерно для таких функций, как сочетание SHIFT+C для копирования). \\ \hline
	Авто-поднятие (Auto-shift) & Т Автоматическое преобразование букв в верхний регистр при длительном удержании определенной клавиши (например, вместо долгого удерживания SHIFT пользователь просто зажимает клавишу и получает заглавную букву). \\ \hline
	RGB-подсветка & Многоцветная светодиодная подсветка клавиатуры, управляемая программно через QMK. Имеет массу предустановленных анимаций и возможность индивидуального программирования освещения каждой клавиши отдельно. \\ \hline
	Перезагрузчик (Bootloader) & Часть прошивки, отвечающая за обновление основного программного обеспечения клавиатуры. Позволяет обновить прошивку без использования специального кабеля или оборудования. \\ \hline
	Однажды-вызванная клавиша (One-Shot Key) & Клавиша, действующая кратковременно при нажатии, освобождая руку пользователя от длительного удерживания модификаторов (например, One-Shot Shift действует, как SHIFT, только на один цикл ввода). \\ \hline
	Microcontroller Unit (MCU) & Микроконтроллер, управляющий внутренней электроникой клавиатуры. Современные платы используют ATMega, Cortex M0/M3 или аналогичные чипы для обеспечения эффективной работы клавиатуры. \\ \hline
	Аппаратное ускорение (Hardware acceleration) & Возможность увеличения быстродействия приложения за счет использования ресурсов видеокарты компьютера. Уменьшает нагрузку на центральный процессор и повышает отзывчивость интерфейса. \\ \hline
	Коллизионная карта (Collision Map) & Таблица, показывающая, какие клавиши конфликтуют друг с другом физически (не могут быть нажаты одновременно из-за конструктивных особенностей клавиатуры). Используется для избежания некорректного поведения клавиш. \\ \hline
	Обратное освещение (Backlight) & Метод считывания состояния клавиш, при котором контроллер последовательно перебирает строки и столбы виртуальной матрицы, фиксируя факт нажатия клавиш. \\ \hline
	Firmware (прошивка) & Специальное программное обеспечение, которое устанавливается на микроконтроллер клавиатуры. Прошивка управляет работой клавиатуры, обрабатывает события нажатия клавиш и передает сигналы на компьютер. \\ \hline
	Keymap (карта клавиш) & Расположение и назначение клавиш на клавиатуре. Карта клавиш хранит информацию о том, какую функцию выполняет каждая отдельная клавиша. В QMK карты клавиш хранятся в виде отдельных файлов (*.c или *.h), что позволяет пользователям самостоятельно настраивать назначение клавиш. \\ \hline
	Layout (макет) & Общая схема распределения клавиш на клавиатуре. Может включать различные типы расстановки клавиш, такие как ортолинейный, сплэт, вогнутый и другие. \\ \hline
	Dynamic Key Overwrite & Функция, позволяющая динамически изменять назначение определенной клавиши в зависимости от ситуации, модификаторов или условий. \\ \hline
	EEPROM (энергонезависимая память) & Небольшой участок памяти, хранящийся на микроконтроллере, в котором хранится конфигурация клавиатуры (картографирование клавиш, параметры подсветки и т.п.). При сбросе прошивки содержимое EEPROM сохраняется. \\ \hline
	Lock Mode (Режим блокировки) & Специальный режим, в котором некоторые клавиши фиксируют свое состояние ("lock"), например, Num Lock или Caps Lock. \\ \hline
	Hotplug Switches (горячая замена свитчей) & Технология, позволяющая менять физические механизмы клавиш (свитчи) без пайки. Часто используется в DIY-клавиатурах. \\ \hline
	Rotary Encoder (энкодер вращения) & Устройство, преобразующее вращательные движения в цифровой сигнал. Часто используется для регулировки громкости, скроллинга или смены слоев. \\ \hline
	Python & Высокоуровневый объектно-ориентированный язык программирования, отличающийся простотой синтаксиса и богатой стандартной библиотекой. Python широко используется для разработки веб-приложений, анализа данных, научных вычислений, автоматизации процессов и создания графических интерфейсов. \\ \hline
	PyQt/PySide & Библиотеки для создания графического интерфейса пользователя. Позволяют строить красивые и современные графические интерфейсы с удобным управлением элементами и событиями. \\ \hline
	hidapi & библиотека для работы с устройствами ввода-вывода (USB, HID), которая позволяет взаимодействовать с физической клавиатурой и считывать/сохранять настройки прошивки. \\ \hline
	json & стандартный модуль Python для работы с форматом JSON, используемый для чтения и записи конфигурационных файлов.
\end{xltabular}
\renewcommand{\arraystretch}{1.0} % восстановление сетки

\subsection{Общие сведения о программно-информационной среде}

Для разработки данного программного продукта был использован язык программирования Python и среда разработки Visual Studio code.

Диаграмма компонентов данного программного продукта представлена на рисунке \ref{componentdiagram:image}

\begin{figure}[H]
	\centering
	\includegraphics[width=0.8\linewidth]{componentdiagram}
	\caption{Диаграмма компонентов}
	\label{componentdiagram:image}
\end{figure}

Данный программный продукт состоит из 16 компонентов:

\begin{enumerate}
\item Работа с файлом. Этот модуль отвечает за работу с файлом.
\item Переопределение ключей. Этот модуль отвечает за переопределение ключей.
\item Безопасность. Этот модуль отвечает за безопасность.
\item Тема. Этот модуль отвечает за тему.
\item Подсветка. Этот модуль отвечает за подсветку.
\item Альтернативные клавиши повтора. Этот модуль отвечает за альтернативные клавиши повтора.
\item Раскладка клавиатуры. Этот модуль отвечает за раскладку клавиатуры.
\item Степ. Этот модуль отвечает за степ.
\item Комбинации. Этот модуль отвечает за комбинации.
\item Слои. Этот модуль отвечает за слои.
\item Языковая раскладка. Этот модуль отвечает за языковую раскладку.
\item Макрос. Этот модуль отвечает за макросы.
\item Настройки прошивки QMK. Этот модуль отвечает за настройки прошивки QMK.
\item Тестирование матрицы клавиатуры. Этот модуль отвечает за тестирование матрицы клавиатуры.
\item Графический интерфейс. Этот модуль отвечает за графический интерфейс.
\item Модуль работы с клавиатурой. Этот модуль отвечает за взаимодействие с клавиатурой.
\end{enumerate}

\subsection{Компоненты системы}

Разработанная система представляет собой графическое приложение для настройки клавиатур с поддержкой прошивок QMK/Vial. Ниже приведено описание основных компонентов (модулей) системы, их файловая структура и связи между ними.

\subsubsection{Модуль обнаружения и управления устройствами}
Файлы: \texttt{autodetect/device\_monitor.py}, \texttt{autodetect/monitor\_thread.py}, \texttt{device\_handler.py}, \texttt{hid\_comm.py}.
Модуль отвечает за периодический поиск HID-устройств с поддержкой расширенного протокола (CustomKeyboard) или устройств в режиме загрузчика (BootloaderDevice). В отдельном потоке \texttt{MonitorThread} выполняется сканирование, фильтрация по magic-строкам в серийном номере и проверка корректности интерфейса. Обнаруженные устройства передаются в главное окно через сигнал \texttt{devices\_updated}. Модуль также обеспечивает открытие/закрытие устройства и подстановку сторонних JSON-определений (sideload).

\subsubsection{Модуль связи с клавиатурой}
Файлы: \texttt{comm/keyboard\_comm.py}, \texttt{comm/base\_comm.py}, \texttt{comm/constants.py}, \texttt{hid\_comm.py}.
Класс \texttt{Keyboard} инкапсулирует всю логику обмена с устройством через HID-отчёты (функция \texttt{hid\_send}). Модуль реализует получение параметров клавиатуры (количество слоёв, матрица, энкодеры), загрузку/сохранение текущей раскладки, управление подсветкой, а также отправку команд динамической конфигурации. Поддерживаются протоколы версий 1–6. \texttt{Keyboard} служит мостом между графическими редакторами и физическим устройством.

\subsubsection{Модуль протоколов динамической конфигурации}
Файлы: \texttt{comm/macro\_protocol.py}, \texttt{comm/tap\_dance\_protocol.py}, \texttt{comm/combo\_protocol.py}, \texttt{comm/key\_override\_protocol.py}, \texttt{comm/alt\_repeat\_key.py}, \texttt{comm/dynamic\_protocol.py}.
Каждый файл реализует отдельный протокол для чтения/записи динамических данных:
\begin{itemize}
	\item макросы (сериализация/десериализация, поддержка задержек и двухбайтовых кодов);
	\item tap dance (до четырёх действий и тайминг);
	\item комбинации (4 клавиши -> результирующая);
	\item переопределения клавиш (условия, модификаторы, слои);
	\item альтернативные клавиши повтора.
\end{itemize}
Все протоколы наследуются от \texttt{BaseProtocol} и используют единый механизм \texttt{\_retrieve\_dynamic\_entries} для пакетного получения данных. Модуль тесно связан с \texttt{Keyboard} (класс-примесь).

\subsubsection{Модуль редактора раскладки клавиатуры}
Файлы: \texttt{config/keymap\_editor.py}, \texttt{config/layout\_editor.py}, \texttt{widgets/keyboard\_widget.py}, \texttt{widgets/key\_widget.py}.
Редактор раскладки (\texttt{KeymapEditor}) отображает графическое представление клавиатуры (виджет \texttt{KeyboardWidget}), позволяет переключать слои, выбирать отдельные клавиши и назначать им произвольные коды (ключарты). При изменении клавиши вызывается метод \texttt{Keyboard.set\_key} для отправки команды на устройство. Редактор параметров раскладки (\texttt{LayoutEditor}) управляет аппаратными опциями (на основе JSON-определения), сохраняет их в \texttt{QSettings} и уведомляет об изменениях через сигнал.

\subsubsection{Модуль макросов}
Файлы: \texttt{config/macro\_recorder.py}, \texttt{macro/macro\_tab.py}, \texttt{macro/macro\_action.py}, \texttt{macro/macro\_optimizer.py}, \texttt{macro/macro\_recorder\_linux.py}.
Позволяет создавать и редактировать макросы в виде последовательности действий (нажатие, отпускание, текст, задержка). Поддерживается интерактивная запись с клавиатуры. Модуль оптимизирует макросы (объединение последовательных нажатий/отпусканий в \texttt{Tap}, замена печатных символов на строки). Данные сериализуются в бинарный буфер и отправляются через \texttt{Keyboard.set\_macro}. Редактор встроен в главное окно как отдельная вкладка.

\subsubsection{Модуль комбинаций (Combos)}
Файлы: \texttt{config/combo\_editor.py}, \texttt{comm/combo\_protocol.py}.
Обеспечивает настройку комбинаций из нескольких одновременно нажатых клавиш (до 4), которые приводят к одному результирующему коду. Интерфейс реализован в виде вкладок, каждая комбинация отображается с помощью виджетов выбора ключей (\texttt{KeyWidget}). Изменения немедленно отправляются на клавиатуру через \texttt{Keyboard.combo\_set}.

\subsubsection{Модуль степов (Tap Dance)}
Файлы: \texttt{config/tap\_dance\_editor.py}, \texttt{comm/tap\_dance\_protocol.py}.
Редактирование действий, назначаемых на один физический переключатель в зависимости от числа нажатий и удержания (одинарный тап, двойной тап, удержание, тап+удержание). Для каждого степа задаётся время тап-термина. Модуль поддерживает отложенную отправку (кнопка Save) и откат (Revert). Протокол использует те же динамические команды, что и комбинации.

\subsubsection{Модуль переопределения клавиш (Key Overrides)}
Файлы: \texttt{config/key\_override\_editor.py}, \texttt{comm/key\_override\_protocol.py}.
Позволяет задать условную замену одной клавиши на другую с учётом активных модификаторов, текущего слоя, а также дополнительных опций (однократное срабатывание, игнорирование отрицательных модификаторов и пр.). Интерфейс включает выбор триггерной клавиши, клавиши-замены, масок модификаторов и слоёв. Каждое переопределение может быть включено/отключено.

\subsubsection{Модуль альтернативных клавиш повтора (Alt Repeat Key)}
Файлы: \texttt{config/alt\_key\_repeat.py}, \texttt{comm/alt\_repeat\_key.py}.
Реализует функцию повторной отправки последней нажатой клавиши с возможностью подмены её на другую (например, \texttt{Alt+Repeat} вместо простого повтора). Настраиваются разрешённые модификаторы и опции двунаправленности. Модуль доступен только при наличии поддержки в прошивке.

\subsubsection{Модуль настройки подсветки}
Файлы: \texttt{config/rgb\_control.py}.
Объединяет три типа подсветки: QMK Backlight, QMK RGBLight и собственную Vial RGB. В зависимости от возможностей устройства отображаются соответствующие элементы управления: яркость, скорость, выбор режима, выбор цвета. Команды отправляются через методы \texttt{Keyboard.set\_qmk\_rgblight\_*} и \texttt{Keyboard.set\_custom\_rgb\_*}. Состояние подсветки можно сохранить командой \texttt{save\_rgb}.

\subsubsection{Модуль настроек QMK}
Файлы: \texttt{config/qmk\_settings\_editor.py}.
Динамически строит интерфейс на основе файла \texttt{qmk\_settings.json} (поставляется с приложением). Поддерживаются параметры типов boolean и integer. Редактор получает список поддерживаемых идентификаторов (QSID) от клавиатуры, загружает текущие значения, позволяет их изменять и сохранять. Использует протокол \texttt{CMD\_QMK\_SETTINGS\_*}.

\subsubsection{Модуль тестирования матрицы}
Файлы: \texttt{config/matrix\_tester.py}.
Отображает виртуальную матрицу клавиш, периодически (таймер 20 мс) опрашивая состояние коммутации через \texttt{Keyboard.matrix\_poll}. Нажатые клавиши подсвечиваются на графическом представлении. Перед тестированием требуется разблокировка устройства (вызов \texttt{Unlocker.unlock}). Модуль работает только при наличии протокола версии 3 и выше.

\subsubsection{Модуль обновления прошивки}
Файлы: \texttt{config/firmware\_updater.py}.
Позволяет прошить клавиатуру файлом с расширением \texttt{.vfw}. При этом происходит проверка целостности (SHA256), сравнение UID устройства с UID прошивки, автоматический переход в режим загрузчика, отправка данных блоками по 64 байта и последующее восстановление сохранённой раскладки (опционально). Модуль взаимодействует с \texttt{BootloaderDevice} и использует низкоуровневую отправку \texttt{pad\_for\_bootloader}.

\subsubsection{Модуль графического интерфейса}
Файлы: \texttt{main\_window.py}, \texttt{widgets/*}, \texttt{keycode\_tabs.py}.
Главное окно (\texttt{MainWindow}) содержит комбобокс выбора устройства, панель вкладок (\texttt{QTabWidget}) для редакторов, меню (файл, безопасность, тема, справка). Виджеты:
\begin{itemize}
	\item \texttt{KeyboardWidget} -- отрисовка клавиатуры по KLE-описанию;
	\item \texttt{KeyWidget} -- компактный выбор клавиши с выпадающей панелью кодов;
	\item \texttt{SquareButton} -- квадратная кнопка с масштабированием;
	\item \texttt{FlowLayout} -- автоматическая упаковка кнопок;
	\item \texttt{KeycodeTabs} -- вкладки с группами клавиш (Basic, ISO/JIS, Layers, Quantum и пр.).
\end{itemize}
Взаимодействие с пользователем осуществляется через сигналы/слоты PyQt.

\subsubsection{Модуль языковых раскладок}
Файлы: \texttt{layouts/*.py}, \texttt{layout\_presets.py}.
Содержит словари переопределения надписей для QWERTY/QWERTZ/AZERTY и альтернативных раскладок (Colemak, Dvorak, Bépo, русская и др.). Выбранная раскладка применяется к отображению клавиш в редакторе (метод \texttt{KeycodeDisplay.set\_keymap\_override}) и сохраняется в \texttt{QSettings}.

\subsubsection{Модуль тем оформления}
Файлы: \texttt{themes.py}.
Определяет набор готовых цветовых схем (Light, Dark, Arc, Nord, Dracula, Catppuccin и пр.). При выборе темы изменяется палитра приложения (\texttt{QPalette}) и стиль Fusion. Настройка сохраняется и применяется после перезапуска.

\subsubsection{Модуль безопасности}
Файлы: \texttt{security\_dialog.py}.
Реализует диалог разблокировки клавиатуры. Показывает графическое изображение клавиш, которые необходимо зажать для перехода в разблокированный режим. В процессе разблокировки отображается прогресс-бар, а при завершении редактирования можно явно заблокировать устройство (метод \texttt{Keyboard.lock}). Модуль используется редакторами перед записью критичных данных (например, макросов или сброса EEPROM).

\subsubsection{Модуль работы с файлами}
Файлы: \texttt{main\_window.py} (методы \texttt{on\_layout\_load}, \texttt{on\_layout\_save}), \texttt{text\_editor.py}.
Обеспечивает загрузку и сохранение файлов конфигурации с расширением \texttt{.vil} (JSON-формат). При сохранении сериализуются: версия, UID клавиатуры, массивы слоёв, опции раскладки, макросы, tap dance, комбинации, переопределения клавиш, настройки QMK. Загрузка восстанавливает состояние через \texttt{Keyboard.restore\_layout}. Также доступен текстовый редактор для ручного редактирования JSON-представления макросов.

\subsection{Связи между компонентами}

На рисунке \ref{componentdiagram:image} (диаграмма компонентов) показано, как модули взаимодействуют друг с другом. Ключевые связи:

\begin{enumerate}
	\item \texttt{MainWindow} агрегирует \texttt{DeviceMonitor} и содержит список редакторов. При выборе устройства вызывается метод \texttt{rebuild} для каждого редактора, передавая ему объект \texttt{CustomKeyboard}.
	\item Редакторы (унаследованные от \texttt{BasicConfig}) через поле \texttt{device} получают доступ к \texttt{Keyboard} и вызывают его методы (\texttt{set\_key}, \texttt{set\_macro}, \texttt{combo\_set}, \texttt{tap\_dance\_set} и т.д.).
	\item \texttt{Keyboard} использует протоколы (\texttt{ProtocolMacro}, \texttt{ProtocolDynamic} и др.) для формирования пакетов и отправки через \texttt{hid\_send}.
	\item Модуль \texttt{device\_handler} создаёт объекты \texttt{CustomKeyboard} или \texttt{BootloaderDevice} на основе HID-дескрипторов, полученных из \texttt{hid\_comm}.
	\item \texttt{KeycodeDisplay} (утилита) глобально переопределяет отображение клавиш в соответствии с выбранной языковой раскладкой и работает со всеми виджетами, зарегистрированными через \texttt{notify\_keymap\_override}.
	\item Система тем (\texttt{themes}) изменяет палитру всего приложения, что влияет на все графические элементы.
	\item Диалог разблокировки (\texttt{Unlocker}) блокирует интерфейс главного окна на время процедуры, чтобы избежать конфликтующих операций.
\end{enumerate}

\subsection{Внешние компоненты системы}

Для работы приложения используются следующие внешние библиотеки и фреймворки:

\begin{enumerate}
	\item \texttt{PyQt5} -- фреймворк для создания графического интерфейса (виджеты, сигналы/слоты, работа с файлами).
	\item \texttt{hidapi} (в зависимости от ОС) -- низкоуровневый доступ к HID-устройствам, пересылка USB-отчётов.
	\item \texttt{fbs} -- утилиты для сборки и определения ресурсов (только для десктопных сборок).
\end{enumerate}

Все внешние компоненты используются через импорты в соответствующих модулях; их версии и лицензии указаны в документации к проекту.

\subsection{Структура файла *.vil}

Файл *.vil -- JSON файл, который включает:

\begin{enumerate}
\item Версию файла (всегда имеет значение 1).
\item Уникальный идентификатор конфигурации.
\item Массив слоёв, где каждый элемент -- массив клавиш.
\item Дополнительные опции раскладки.
\item Массив макросов, где каждый макрос -- это массив действий.
\item Массив степов, где каждый степ -- массив из 5 элементов (при одинарном нажатии, при двойном нажатии, при удержании, при нажатии и удержании, задержка между действиями).
\item Массив комбо, где каждое комбо -- массив из 5 элементов (4 клавиши и результирующая клавиша).
\item Массив переопределения ключей, где каждое переопределение -- объект, в котором есть следующие поля:

\begin{itemize}
	\item нажимаемая клавиша;
	\item результирующая клавиша;
	\item в каких слоях используется;
	\item маски модификаторов;
	\item дополнительные флаги.
\end{itemize}

Пример файла *.vil представлен на рисунке \ref{vilfile:image}:

\begin{figure}[H]
	\lstinputlisting[language={}, frame=none]{config.vil}
	\caption{Пример файла *.vil}
	\label{vilfile:image}
\end{figure}

\item Массив альтернативных клавиш повтора, где каждый элемент -- объект, в котором есть следующие поля:

\begin{itemize}
	\item нажимаемая клавиша;
	\item альтернативная клавиша;
	\item разрешённый модификаторы;
	\item дополнительные флаги.
\end{itemize}

\item Объект настроек прошивки QMK, который содержит все настройки прошивки QMK.
\end{enumerate}

\subsection{Проектирование пользовательского интерфейса программы}

На рисунке \ref{mainlayout:image} изображена главная страница программы.

\begin{figure}[H]
	\centering
	\includegraphics[width=1\linewidth]{mainlayout}
	\caption{Интерфейс главной страницы}
	\label{mainlayout:image}
\end{figure}

В таблице \ref{main:layouttable} представлено описание объектов данного интерфейса.

\renewcommand{\arraystretch}{0.8} % уменьшение расстояний до сетки таблицы
\begin{xltabular}{\textwidth}{|p{2.25cm}|p{2.25cm}|X|}
	\caption{Описание объектов интерфейса главной страницы\label{main:layouttable}}\\
	\hline \centrow \setlength{\baselineskip}{0.7\baselineskip} Номер объекта & \centrow \setlength{\baselineskip}{0.7\baselineskip} Тип объекта & \centrow \setlength{\baselineskip}{0.7\baselineskip} Действие/описание объекта \\
	\hline \centrow 1 & \centrow 2 & \centrow 3\\ \hline
	\endfirsthead
	\caption*{Продолжение таблицы \ref{main:layouttable}}\\
	\hline \centrow 1 & \centrow 2 & \centrow 3\\ \hline
	\finishhead
	1 & Кнопка & Открывает меню "<File"> \\ \hline
	2 & Кнопка & Открывает меню "<Keyboard layout"> \\ \hline
	3 & Кнопка & Открывает меню "<Security"> \\ \hline
	4 & Кнопка & Открывает меню "<Theme"> \\ \hline
	5 & Кнопка & Открывает меню "<About"> \\ \hline
	6 & Кнопка & Открывает меню "<Keymap"> \\ \hline
	7 & Кнопка & Открывает меню "<Macros"> \\ \hline
	8 & Кнопка & Открывает меню "<Lightning"> \\ \hline
	9 & Кнопка & Открывает меню "<Tap Dance"> \\ \hline
	10 & Кнопка & Открывает меню "<Combos"> \\ \hline
	11 & Кнопка & Открывает меню "<Key Overrides"> \\ \hline
	12 & Кнопка & Открывает меню "<Alt Repeat Key"> \\ \hline
	13 & Кнопка & Открывает меню "<QMK Settings"> \\ \hline
	14 & Кнопки & Позволяют выбрать слой клавиш от 1 до 14 \\ \hline
	15 & Кнопки & Позволяют настроить каждый слой клавиатуры от 1 до 14 \\ \hline
	16 & Кнопка & Открывает базовые клавиши для их использования \\ \hline
	17 & Кнопка & Открывает клавиши для Европейской и Японской клавиатуры \\ \hline
	18 & Кнопка & Открывает кнопки слоёв \\ \hline
	19 & Кнопка & Открывает клавиши QMK прошивки \\ \hline
	20 & Кнопка & Открывает клавиши подсветки \\ \hline
	21 & Кнопка & Открывает клавиши приложений, медиа и мыши \\ \hline
	22 & Кнопка & Открывает клавиши Tap Dance \\ \hline
	23 & Кнопка & Открывает пользовательские клавиши \\ \hline
	24 & Кнопка & Открывает клавиши макросов \\ \hline
	25 & Кнопка & Поле выбора нужной кнопки для настройки слоя клавиатуры \\ \hline
	26 & Кнопка & Сбрасывает все настройки до первоначальных
\end{xltabular}
\renewcommand{\arraystretch}{1.0} % восстановление сетки

На рисунке \ref{file:layoutimage} изображено выпадающее окно "<File">.

\begin{figure}[H]
	\centering
	\includegraphics[width=0.6\linewidth]{filelayout}
	\caption{Выпадающее окно "<File">}
	\label{file:layoutimage}
\end{figure}

В таблице \ref{file:layouttable} Выпадающее окно "<File">.

\renewcommand{\arraystretch}{0.8} % уменьшение расстояний до сетки таблицы
\begin{xltabular}{\textwidth}{|p{2.25cm}|p{2.25cm}|X|}
	\caption{Описание объектов выпадающего окна "<File">\label{file:layouttable}}\\
	\hline \centrow \setlength{\baselineskip}{0.7\baselineskip} Номер объекта & \centrow \setlength{\baselineskip}{0.7\baselineskip} Тип объекта & \centrow \setlength{\baselineskip}{0.7\baselineskip} Действие/описание объекта \\
	\hline \centrow 1 & \centrow 2 & \centrow 3\\ \hline
	\endfirsthead
	\caption*{Продолжение таблицы \ref{file:layouttable}}\\
	\hline \centrow 1 & \centrow 2 & \centrow 3\\ \hline
	\finishhead
	1 & Кнопка & Загружает файл с ПК \\ \hline
	2 & Кнопка & Сохраняет текущую раскладку в файл на ПК \\ \hline
	3 & Кнопка & Дополнительная загрузка через JSON \\ \hline
	4 & Кнопка & Скачать VIA \\ \hline
	5 & Кнопка & Загружает фиктивный JSON \\ \hline
	6 & Кнопка & Закрывает выпадающее окно
\end{xltabular}
\renewcommand{\arraystretch}{1.0} % восстановление сетки

На рисунке \ref{Keyboard layout:layoutimage} изображена выпадающее окно "<Keyboard layout">.

\begin{figure}[H]
	\centering
	\includegraphics[width=0.4\linewidth]{Keyboard layout layout}
	\caption{Выпадающее окно "<Keyboard layout">}
	\label{Keyboard layout:layoutimage}
\end{figure}

В таблице \ref{Keyboard layout:layouttable} представлено описание объектов данного выпадающего окна.

\renewcommand{\arraystretch}{0.8} % уменьшение расстояний до сетки таблицы
\begin{xltabular}{\textwidth}{|p{2.25cm}|p{2.25cm}|X|}
	\caption{Описание объектов интерфейса выпадающего окна "<Keyboard layout">\label{Keyboard layout:layouttable}}\\
	\hline \centrow \setlength{\baselineskip}{0.7\baselineskip} Номер объекта & \centrow \setlength{\baselineskip}{0.7\baselineskip} Тип объекта & \centrow \setlength{\baselineskip}{0.7\baselineskip} Действие/описание объекта \\
	\hline \centrow 1 & \centrow 2 & \centrow 3\\ \hline
	\endfirsthead
	\caption*{Продолжение таблицы \ref{Keyboard layout:layouttable}}\\
	\hline \centrow 1 & \centrow 2 & \centrow 3\\ \hline
	\finishhead
	1 & Флажки & Выбор нужной языковой раскладки клавиатуры
\end{xltabular}
\renewcommand{\arraystretch}{1.0} % восстановление сетки

На рисунке \ref{Security:layoutimage} изображено выпадающее окно "<Security">.

\begin{figure}[H]
	\centering
	\includegraphics[width=0.6\linewidth]{Securitylayout}
	\caption{выпадающее окно "<Security">}
	\label{Security:layoutimage}
\end{figure}

В таблице \ref{Security:layouttable} представлено описание объектов данного выпадающего окна.

\renewcommand{\arraystretch}{0.8} % уменьшение расстояний до сетки таблицы
\begin{xltabular}{\textwidth}{|p{2.25cm}|p{2.25cm}|X|}
	\caption{Описание объектов интерфейса выпадающего окна "<Security">\label{Security:layouttable}}\\
	\hline \centrow \setlength{\baselineskip}{0.7\baselineskip} Номер объекта & \centrow \setlength{\baselineskip}{0.7\baselineskip} Тип объекта & \centrow \setlength{\baselineskip}{0.7\baselineskip} Действие/описание объекта \\
	\hline \centrow 1 & \centrow 2 & \centrow 3\\ \hline
	\endfirsthead
	\caption*{Продолжение таблицы \ref{Security:layouttable}}\\
	\hline \centrow 1 & \centrow 2 & \centrow 3\\ \hline
	\finishhead
	1 & Кнопка & Разблокировка клавиатуры для изменения в ней параметров \\ \hline
	2 & Кнопка & Блокировка клавиатуры для невозможности измены её параметров \\ \hline
	3 & Кнопка & Перезагрузка клавиатуры
\end{xltabular}
\renewcommand{\arraystretch}{1.0} % восстановление сетки

На рисунке \ref{Theme:layoutimage} изображено выпадающее окно "<Theme">.

\begin{figure}[H]
	\centering
	\includegraphics[width=0.4\linewidth]{Themelayout}
	\caption{выпадающее окно "<Theme">}
	\label{Theme:layoutimage}
\end{figure}

В таблице \ref{Theme:layouttable} представлено описание объектов данного выпадающего окна.

\renewcommand{\arraystretch}{0.8} % уменьшение расстояний до сетки таблицы
\begin{xltabular}{\textwidth}{|p{2.25cm}|p{2.25cm}|X|}
	\caption{Описание объектов интерфейса выпадающего окна "<Theme">\label{Theme:layouttable}}\\
	\hline \centrow \setlength{\baselineskip}{0.7\baselineskip} Номер объекта & \centrow \setlength{\baselineskip}{0.7\baselineskip} Тип объекта & \centrow \setlength{\baselineskip}{0.7\baselineskip} Действие/описание объекта \\
	\hline \centrow 1 & \centrow 2 & \centrow 3\\ \hline
	\endfirsthead
	\caption*{Продолжение таблицы \ref{Theme:layouttable}}\\
	\hline \centrow 1 & \centrow 2 & \centrow 3\\ \hline
	\finishhead
	1 & Флажки & Выбор нужной цветовой темы программы
\end{xltabular}
\renewcommand{\arraystretch}{1.0} % восстановление сетки

На рисунке \ref{About:layoutimage} изображена страница "<About">.

\begin{figure}[H]
	\centering
	\includegraphics[width=0.6\linewidth]{Aboutlayout}
	\caption{страница "<About">}
	\label{About:layoutimage}
\end{figure}

В таблице \ref{About:layouttable} представлено описание объектов данного интерфейса.

\renewcommand{\arraystretch}{0.8} % уменьшение расстояний до сетки таблицы
\begin{xltabular}{\textwidth}{|p{2.25cm}|p{2.25cm}|X|}
	\caption{Описание объектов интерфейса страницы "<About">\label{About:layouttable}}\\
	\hline \centrow \setlength{\baselineskip}{0.7\baselineskip} Номер объекта & \centrow \setlength{\baselineskip}{0.7\baselineskip} Тип объекта & \centrow \setlength{\baselineskip}{0.7\baselineskip} Действие/описание объекта \\
	\hline \centrow 1 & \centrow 2 & \centrow 3\\ \hline
	\endfirsthead
	\caption*{Продолжение таблицы \ref{About:layouttable}}\\
	\hline \centrow 1 & \centrow 2 & \centrow 3\\ \hline
	\finishhead
	1 & Кнопка & Выводит всю информацию о клавиатуре \\ \hline
	2 & Кнопка & Выводит краткую информацию о программе
\end{xltabular}
\renewcommand{\arraystretch}{1.0} % восстановление сетки

На рисунке \ref{Macros:layoutimage} изображена страница "<Macros">.

\begin{figure}[H]
	\centering
	\includegraphics[width=1\linewidth]{Macroslayout}
	\caption{Страница "<Macros">}
	\label{Macros:layoutimage}
\end{figure}

В таблице \ref{Macros:layouttable} представлено описание объектов данного интерфейса.

\renewcommand{\arraystretch}{0.8} % уменьшение расстояний до сетки таблицы
\begin{xltabular}{\textwidth}{|p{2.25cm}|p{2.25cm}|X|}
	\caption{Описание объектов интерфейса страницы "<Macros">\label{Macros:layouttable}}\\
	\hline \centrow \setlength{\baselineskip}{0.7\baselineskip} Номер объекта & \centrow \setlength{\baselineskip}{0.7\baselineskip} Тип объекта & \centrow \setlength{\baselineskip}{0.7\baselineskip} Действие/описание объекта \\
	\hline \centrow 1 & \centrow 2 & \centrow 3\\ \hline
	\endfirsthead
	\caption*{Продолжение таблицы \ref{Macros:layouttable}}\\
	\hline \centrow 1 & \centrow 2 & \centrow 3\\ \hline
	\finishhead
	1 & Кнопки & Кнопки для выбора профиля макроса \\ \hline
	2 & Кнопки & Для удобного пролистывания профилей макросов \\ \hline
	3 & Кнопки & Удаляют действия макроса \\ \hline
	4 & Кнопки & Перемещают действия вверх или вниз \\ \hline
	5 & Кнопки & Открывают панель выбора кнопки из той или иной раскладки клавиатуры \\ \hline
	6 & Числовое поле & Позволяет ввести в него то количество мили секунд, которое будет выполняться то или иное действие \\ \hline
	7 & Радио кнопки & На выбор несколько вариантов действий макроса \\ \hline
	8 & Кнопка & Открывает текстовый редактор \\ \hline
	9 & Кнопка & Добавляет одно действие для макроса \\ \hline
	10 & Кнопка & Добавляет в макрос действие нажатия кнопки Enter \\ \hline
	11 & Радио кнопки & Позволяет сохранить макрос, либо добавить его к уже существующему макросу \\ \hline
	12 & Кнопка & Сохраняет макрос \\ \hline
	13 & Кнопка & Удаляет все действия с макросом, очищая рабочую зону \\ \hline
	14 & Label & Показывает какое количество памяти занимает макрос
\end{xltabular}
\renewcommand{\arraystretch}{1.0} % восстановление сетки

На рисунке \ref{Lightning:layoutimage} изображена страница "<Lightning">.

\begin{figure}[H]
	\centering
	\includegraphics[width=1\linewidth]{Lightninglayout}
	\caption{Страница "<Lightning">}
	\label{Lightning:layoutimage}
\end{figure}

В таблице \ref{Lightning:layouttable} представлено описание объектов данного интерфейса.

\renewcommand{\arraystretch}{0.8} % уменьшение расстояний до сетки таблицы
\begin{xltabular}{\textwidth}{|p{2.25cm}|p{2.25cm}|X|}
	\caption{Описание объектов интерфейса страницы "<Lightning">\label{Lightning:layouttable}}\\
	\hline \centrow \setlength{\baselineskip}{0.7\baselineskip} Номер объекта & \centrow \setlength{\baselineskip}{0.7\baselineskip} Тип объекта & \centrow \setlength{\baselineskip}{0.7\baselineskip} Действие/описание объекта \\
	\hline \centrow 1 & \centrow 2 & \centrow 3\\ \hline
	\endfirsthead
	\caption*{Продолжение таблицы \ref{Lightning:layouttable}}\\
	\hline \centrow 1 & \centrow 2 & \centrow 3\\ \hline
	\finishhead
	1 & Радио кнопки & Позволяет выбрать один из многих типов подсветки клавиатуры \\ \hline
	2 & Кнопка & Открывает меню точной настройки цвета по трём каналам RGB \\ \hline
	3 & Ползунок & Позволяет настроить яркость подсветки \\ \hline
	4 & Ползунок & Позволяет настроить скорость смены подсветки \\ \hline
	5 & Кнопка & Сохраняет подсветку \\ \hline
	6 & Подписи & Непосредственно подписывают то, к чему относятся элементы возле них
\end{xltabular}
\renewcommand{\arraystretch}{1.0} % восстановление сетки

На рисунке \ref{Tap Dance:layoutimage} изображена страница "<Tap Dance">.

\begin{figure}[H]
	\centering
	\includegraphics[width=1\linewidth]{Tap Dancelayout}
	\caption{Страница "<Tap Dance">}
	\label{Tap Dance:layoutimage}
\end{figure}

В таблице \ref{Tap Dance:layouttable} представлено описание объектов данного интерфейса.

\renewcommand{\arraystretch}{0.8} % уменьшение расстояний до сетки таблицы
\begin{xltabular}{\textwidth}{|p{2.25cm}|p{2.25cm}|X|}
	\caption{Описание объектов интерфейса страницы "<Tap Dance">\label{Tap Dance:layouttable}}\\
	\hline \centrow \setlength{\baselineskip}{0.7\baselineskip} Номер объекта & \centrow \setlength{\baselineskip}{0.7\baselineskip} Тип объекта & \centrow \setlength{\baselineskip}{0.7\baselineskip} Действие/описание объекта \\
	\hline \centrow 1 & \centrow 2 & \centrow 3\\ \hline
	\endfirsthead
	\caption*{Продолжение таблицы \ref{Tap Dance:layouttable}}\\
	\hline \centrow 1 & \centrow 2 & \centrow 3\\ \hline
	\finishhead
	1 & Кнопки & Кнопки выбора одного из "<Tap Dance"> \\ \hline
	2 & Label & Отображают информацию о элементах справа \\ \hline
	3 & Label & Описание поля для ввода данных \\ \hline
	4 & Кнопки & Открывают панель выбора клавиш раскладки клавиатуры \\ \hline
	5 & Числовое поле & Позволяет ввести в него то количество мили секунд, которое будет дано на выполнение "<Tap Dance"> \\ \hline
	6 & Кнопка & Сохраняет готовый "<Tap Dance"> \\ \hline
	7 & Кнопка & Очищает рабочее поле
\end{xltabular}
\renewcommand{\arraystretch}{1.0} % восстановление сетки

На рисунке \ref{Combos:layoutimage} изображена страница "<Combos">.

\begin{figure}[H]
	\centering
	\includegraphics[width=1\linewidth]{Comboslayout}
	\caption{Страница "<Combos">}
	\label{Combos:layoutimage}
\end{figure}

В таблице \ref{Combos:layouttable} представлено описание объектов данного интерфейса.

\renewcommand{\arraystretch}{0.8} % уменьшение расстояний до сетки таблицы
\begin{xltabular}{\textwidth}{|p{2.25cm}|p{2.25cm}|X|}
	\caption{Описание объектов интерфейса страницы "<Combos">\label{Combos:layouttable}}\\
	\hline \centrow \setlength{\baselineskip}{0.7\baselineskip} Номер объекта & \centrow \setlength{\baselineskip}{0.7\baselineskip} Тип объекта & \centrow \setlength{\baselineskip}{0.7\baselineskip} Действие/описание объекта \\
	\hline \centrow 1 & \centrow 2 & \centrow 3\\ \hline
	\endfirsthead
	\caption*{Продолжение таблицы \ref{Combos:layouttable}}\\
	\hline \centrow 1 & \centrow 2 & \centrow 3\\ \hline
	\finishhead
	1 & Кнопки & Кнопки выбора одного из "<Combos"> \\ \hline
	2 & Label & Отображают информацию о элементах справа \\ \hline
	3 & Кнопки & Открывают панель выбора клавиш раскладки клавиатуры
\end{xltabular}
\renewcommand{\arraystretch}{1.0} % восстановление сетки

На рисунке \ref{Key Overrides:layoutimage} изображена страница "<Key Overrides">.

\begin{figure}[H]
	\centering
	\includegraphics[width=1\linewidth]{Key Overrideslayout}
	\caption{Страница "<Key Overrides">}
	\label{Key Overrides:layoutimage}
\end{figure}

В таблице \ref{Key Overrides:layouttable} представлено описание объектов данного интерфейса.

\renewcommand{\arraystretch}{0.8} % уменьшение расстояний до сетки таблицы
\begin{xltabular}{\textwidth}{|p{2.25cm}|p{2.25cm}|X|}
	\caption{Описание объектов интерфейса страницы "<Key Overrides">\label{Key Overrides:layouttable}}\\
	\hline \centrow \setlength{\baselineskip}{0.7\baselineskip} Номер объекта & \centrow \setlength{\baselineskip}{0.7\baselineskip} Тип объекта & \centrow \setlength{\baselineskip}{0.7\baselineskip} Действие/описание объекта \\
	\hline \centrow 1 & \centrow 2 & \centrow 3\\ \hline
	\endfirsthead
	\caption*{Продолжение таблицы \ref{Key Overrides:layouttable}}\\
	\hline \centrow 1 & \centrow 2 & \centrow 3\\ \hline
	\finishhead
	1 & Кнопки & Кнопки выбора одного из "<Key Overrides"> \\ \hline
	2 & Label & Отображают информацию о элементах справа \\ \hline
	3 & Флажки & Позволяют настроить нужные слои, которые будут активны \\ \hline
	4 & Кнопки & Позволяют активировать/выключить все флажки для настройки активных слоёв \\ \hline
	5 & Кнопка & При нажатии задаётся кнопка триггера из всех раскладок клавиатур \\ \hline
	6 & Флажки & Выбирают активные функции \\ \hline
	7 & Кнопка & При нажатии задаётся кнопка, на которую будет заменяться предыдущая кнопка \\ \hline
	8 & Label & Показывает информацию для флажков слева \\ \hline
	9 & Флажок & При нажатии включается та или иная настройка
\end{xltabular}
\renewcommand{\arraystretch}{1.0} % восстановление сетки

На рисунке \ref{Alt Repeat Key:layoutimage} изображена страница "<Alt Repeat Key">.

\begin{figure}[H]
	\centering
	\includegraphics[width=1\linewidth]{Alt Repeat Keylayout}
	\caption{Страница "<Alt Repeat Key">}
	\label{Alt Repeat Key:layoutimage}
\end{figure}

В таблице \ref{Alt Repeat Key:layouttable} представлено описание объектов данного интерфейса.

\renewcommand{\arraystretch}{0.8} % уменьшение расстояний до сетки таблицы
\begin{xltabular}{\textwidth}{|p{2.25cm}|p{2.25cm}|X|}
	\caption{Описание объектов интерфейса страницы "<Alt Repeat Key">\label{Alt Repeat Key:layouttable}}\\
	\hline \centrow \setlength{\baselineskip}{0.7\baselineskip} Номер объекта & \centrow \setlength{\baselineskip}{0.7\baselineskip} Тип объекта & \centrow \setlength{\baselineskip}{0.7\baselineskip} Действие/описание объекта \\
	\hline \centrow 1 & \centrow 2 & \centrow 3\\ \hline
	\endfirsthead
	\caption*{Продолжение таблицы \ref{Alt Repeat Key:layouttable}}\\
	\hline \centrow 1 & \centrow 2 & \centrow 3\\ \hline
	\finishhead
	1 & Кнопки & Кнопки выбора одного из "<Alt Repeat Key"> \\ \hline
	2 & Label & Отображают информацию о элементах справа \\ \hline
	3 & Флажок & Позволяет настроить включённость настроек \\ \hline
	4 & Кнопка & При нажатии можно выбрать кнопку, которая будет заменена на альтернативную, и предыдущую кнопку \\ \hline
	5 & Флажки & Позволяют настроить "<Alt Repeat Key">
\end{xltabular}
\renewcommand{\arraystretch}{1.0} % восстановление сетки

На рисунке \ref{QMK Settings:layoutimage} изображена страница "<QMK Settings">.

\begin{figure}[H]
	\centering
	\includegraphics[width=1\linewidth]{QMK Settingslayout}
	\caption{Страница "<QMK Settings">}
	\label{QMK Settings:layoutimage}
\end{figure}

В таблице \ref{QMK Settings:layouttable} представлено описание объектов данного интерфейса.

\renewcommand{\arraystretch}{0.8} % уменьшение расстояний до сетки таблицы
\begin{xltabular}{\textwidth}{|p{2.25cm}|p{2.25cm}|X|}
	\caption{Описание объектов интерфейса страницы "<QMK Settings">\label{QMK Settings:layouttable}}\\
	\hline \centrow \setlength{\baselineskip}{0.7\baselineskip} Номер объекта & \centrow \setlength{\baselineskip}{0.7\baselineskip} Тип объекта & \centrow \setlength{\baselineskip}{0.7\baselineskip} Действие/описание объекта \\
	\hline \centrow 1 & \centrow 2 & \centrow 3\\ \hline
	\endfirsthead
	\caption*{Продолжение таблицы \ref{QMK Settings:layouttable}}\\
	\hline \centrow 1 & \centrow 2 & \centrow 3\\ \hline
	\finishhead
	1 & Кнопки & Кнопки выбора одного из "<QMK Settings"> \\ \hline
	2 & Label & Отображают информацию о элементах справа \\ \hline
	3 & Флажки & Позволяет настроить включённость настроек \\ \hline
	4 & Кнопка & Сбрасывает настройки \\ \hline
	5 & Кнопка & Возвращает предыдущие настройки \\ \hline
	6 & Кнопка & Сохраняет настройки
\end{xltabular}
\renewcommand{\arraystretch}{1.0} % восстановление сетки

На рисунке \ref{Grave Escape:layoutimage} изображена страница "<Grave Escape">.

\begin{figure}[H]
	\centering
	\includegraphics[width=1\linewidth]{Grave Escapelayout}
	\caption{Страница "<Grave Escape">}
	\label{Grave Escape:layoutimage}
\end{figure}

В таблице \ref{Grave Escape:layouttable} представлено описание объектов данного интерфейса.

\renewcommand{\arraystretch}{0.8} % уменьшение расстояний до сетки таблицы
\begin{xltabular}{\textwidth}{|p{2.25cm}|p{2.25cm}|X|}
	\caption{Описание объектов интерфейса страницы "<Grave Escape">\label{Grave Escape:layouttable}}\\
	\hline \centrow \setlength{\baselineskip}{0.7\baselineskip} Номер объекта & \centrow \setlength{\baselineskip}{0.7\baselineskip} Тип объекта & \centrow \setlength{\baselineskip}{0.7\baselineskip} Действие/описание объекта \\
	\hline \centrow 1 & \centrow 2 & \centrow 3\\ \hline
	\endfirsthead
	\caption*{Продолжение таблицы \ref{Grave Escape:layouttable}}\\
	\hline \centrow 1 & \centrow 2 & \centrow 3\\ \hline
	\finishhead
	1 & Label & Отображают информацию о элементах справа \\ \hline
	2 & Флажки & Позволяет настроить включённость настроек \\ \hline
	3 & Кнопки & Сохранить, вернуть, сбросить
\end{xltabular}
\renewcommand{\arraystretch}{1.0} % восстановление сетки

На рисунке \ref{Tap-Hold:layoutimage} изображена страница "<Tap-Hold">.

\begin{figure}[H]
	\centering
	\includegraphics[width=1\linewidth]{Tap-Holdlayout}
	\caption{Страница "<Tap-Hold">}
	\label{Tap-Hold:layoutimage}
\end{figure}

В таблице \ref{Tap-Hold:layouttable} представлено описание объектов данного интерфейса.

\renewcommand{\arraystretch}{0.8} % уменьшение расстояний до сетки таблицы
\begin{xltabular}{\textwidth}{|p{2.25cm}|p{2.25cm}|X|}
	\caption{Описание объектов интерфейса страницы "<Tap-Hold">\label{Tap-Hold:layouttable}}\\
	\hline \centrow \setlength{\baselineskip}{0.7\baselineskip} Номер объекта & \centrow \setlength{\baselineskip}{0.7\baselineskip} Тип объекта & \centrow \setlength{\baselineskip}{0.7\baselineskip} Действие/описание объекта \\
	\hline \centrow 1 & \centrow 2 & \centrow 3\\ \hline
	\endfirsthead
	\caption*{Продолжение таблицы \ref{Tap-Hold:layouttable}}\\
	\hline \centrow 1 & \centrow 2 & \centrow 3\\ \hline
	\finishhead
	1 & Label & Отображают информацию о элементах справа \\ \hline
	2 & Числовое поле & Вносит числовые характеристики \\ \hline
	3 & Флажки & Позволяют настроить параметры, описанные слева \\ \hline
	4 & Числовые поля & Позволяют ввести числа для более точной настройки \\ \hline
	5 & Флажок & Настройка удерживания хордовых \\ \hline
	6 & Числовое поле & Позволяет ввести числовые переменные для более точной настройки \\ \hline
	7 & Кнопки & Сохранить, вернуть, сбросить
\end{xltabular}
\renewcommand{\arraystretch}{1.0} % восстановление сетки

На рисунке \ref{Auto Shift:layoutimage} изображена страница "<Auto Shift">.

\begin{figure}[H]
	\centering
	\includegraphics[width=1\linewidth]{Auto Shiftlayout}
	\caption{Страница "<Auto Shift">}
	\label{Auto Shift:layoutimage}
\end{figure}

В таблице \ref{Auto Shift:layouttable} представлено описание объектов данного интерфейса.

\renewcommand{\arraystretch}{0.8} % уменьшение расстояний до сетки таблицы
\begin{xltabular}{\textwidth}{|p{2.25cm}|p{2.25cm}|X|}
	\caption{Описание объектов интерфейса страницы "<Auto Shift">\label{Auto Shift:layouttable}}\\
	\hline \centrow \setlength{\baselineskip}{0.7\baselineskip} Номер объекта & \centrow \setlength{\baselineskip}{0.7\baselineskip} Тип объекта & \centrow \setlength{\baselineskip}{0.7\baselineskip} Действие/описание объекта \\
	\hline \centrow 1 & \centrow 2 & \centrow 3\\ \hline
	\endfirsthead
	\caption*{Продолжение таблицы \ref{Auto Shift:layouttable}}\\
	\hline \centrow 1 & \centrow 2 & \centrow 3\\ \hline
	\finishhead
	1 & Label & Отображают информацию о элементах справа \\ \hline
	2 & Флажки & Включают, включают для модификаторов "<Auto Shift"> \\ \hline
	3 & Числовое поле & Позволяет настроить "<Timeout"> \\ \hline
	4 & Флажки & Позволяют настроить параметры, описанные слева \\ \hline
	5 & Кнопки & Сохранить, вернуть, сбросить
\end{xltabular}
\renewcommand{\arraystretch}{1.0} % восстановление сетки

На рисунке \ref{Combo:layoutimage} изображена страница "<Combo">.

\begin{figure}[H]
	\centering
	\includegraphics[width=1\linewidth]{Combolayout}
	\caption{Страница "<Combo">}
	\label{Combo:layoutimage}
\end{figure}

В таблице \ref{Combo:layouttable} представлено описание объектов данного интерфейса.

\renewcommand{\arraystretch}{0.8} % уменьшение расстояний до сетки таблицы
\begin{xltabular}{\textwidth}{|p{2.25cm}|p{2.25cm}|X|}
	\caption{Описание объектов интерфейса страницы "<Combo">\label{Combo:layouttable}}\\
	\hline \centrow \setlength{\baselineskip}{0.7\baselineskip} Номер объекта & \centrow \setlength{\baselineskip}{0.7\baselineskip} Тип объекта & \centrow \setlength{\baselineskip}{0.7\baselineskip} Действие/описание объекта \\
	\hline \centrow 1 & \centrow 2 & \centrow 3\\ \hline
	\endfirsthead
	\caption*{Продолжение таблицы \ref{Combo:layouttable}}\\
	\hline \centrow 1 & \centrow 2 & \centrow 3\\ \hline
	\finishhead
	1 & Label & Отображают информацию о элементах справа \\ \hline
	2 & Числовое поле & Позволяет настроить "<Timeout"> между комбинациями \\ \hline
	3 & Кнопки & Сохранить, вернуть, сбросить
\end{xltabular}
\renewcommand{\arraystretch}{1.0} % восстановление сетки

На рисунке \ref{One Shot Keys:layoutimage} изображена страница "<One Shot Keys">.

\begin{figure}[H]
	\centering
	\includegraphics[width=1\linewidth]{One Shot Keyslayout}
	\caption{Страница "<One Shot Keys">}
	\label{One Shot Keys:layoutimage}
\end{figure}

В таблице \ref{One Shot Keys:layouttable} представлено описание объектов данного интерфейса.

\renewcommand{\arraystretch}{0.8} % уменьшение расстояний до сетки таблицы
\begin{xltabular}{\textwidth}{|p{2.25cm}|p{2.25cm}|X|}
	\caption{Описание объектов интерфейса страницы "<One Shot Keys">\label{One Shot Keys:layouttable}}\\
	\hline \centrow \setlength{\baselineskip}{0.7\baselineskip} Номер объекта & \centrow \setlength{\baselineskip}{0.7\baselineskip} Тип объекта & \centrow \setlength{\baselineskip}{0.7\baselineskip} Действие/описание объекта \\
	\hline \centrow 1 & \centrow 2 & \centrow 3\\ \hline
	\endfirsthead
	\caption*{Продолжение таблицы \ref{One Shot Keys:layouttable}}\\
	\hline \centrow 1 & \centrow 2 & \centrow 3\\ \hline
	\finishhead
	1 & Label & Отображают информацию о элементах справа \\ \hline
	2 & Числовые поля & Позволяют настроить количество раз, которое клавиша удерживается для повторного нажатия, время в мили секундах до отпускания клавиши \\ \hline
	3 & Кнопки & Сохранить, вернуть, сбросить
\end{xltabular}
\renewcommand{\arraystretch}{1.0} % восстановление сетки

На рисунке \ref{Matrix tester:layoutimage} изображена страница "<Matrix tester">.

\begin{figure}[H]
	\centering
	\includegraphics[width=1\linewidth]{Matrix testerlayout}
	\caption{Страница "<Matrix tester">}
	\label{Matrix tester:layoutimage}
\end{figure}

В таблице \ref{Matrix tester:layouttable} представлено описание объектов данного интерфейса.

\renewcommand{\arraystretch}{0.8} % уменьшение расстояний до сетки таблицы
\begin{xltabular}{\textwidth}{|p{2.25cm}|p{2.25cm}|X|}
	\caption{Описание объектов интерфейса страницы "<Matrix tester">\label{Matrix tester:layouttable}}\\
	\hline \centrow \setlength{\baselineskip}{0.7\baselineskip} Номер объекта & \centrow \setlength{\baselineskip}{0.7\baselineskip} Тип объекта & \centrow \setlength{\baselineskip}{0.7\baselineskip} Действие/описание объекта \\
	\hline \centrow 1 & \centrow 2 & \centrow 3\\ \hline
	\endfirsthead
	\caption*{Продолжение таблицы \ref{Matrix tester:layouttable}}\\
	\hline \centrow 1 & \centrow 2 & \centrow 3\\ \hline
	\finishhead
	1 & Кнопки & Каждая из нажатых кнопок на клавиатуре отображается на данной матрице кнопок \\ \hline
	2 & Кнопка & Обнуляет матрицу
\end{xltabular}
\renewcommand{\arraystretch}{1.0} % восстановление сетки

На рисунке \ref{basekeymap:layoutimage} изображена страница базовой раскладки клавиатуры.

\begin{figure}[H]
	\centering
	\includegraphics[width=1\linewidth]{basekeymaplayout}
	\caption{Страница базовой раскладки клавиатуры}
	\label{basekeymap:layoutimage}
\end{figure}

В таблице \ref{basekeymap:layouttable} представлено описание объектов данного интерфейса.

\renewcommand{\arraystretch}{0.8} % уменьшение расстояний до сетки таблицы
\begin{xltabular}{\textwidth}{|p{2.25cm}|p{2.25cm}|X|}
	\caption{Описание объектов интерфейса страницы базовой раскладки клавиатуры\label{basekeymap:layouttable}}\\
	\hline \centrow \setlength{\baselineskip}{0.7\baselineskip} Номер объекта & \centrow \setlength{\baselineskip}{0.7\baselineskip} Тип объекта & \centrow \setlength{\baselineskip}{0.7\baselineskip} Действие/описание объекта \\
	\hline \centrow 1 & \centrow 2 & \centrow 3\\ \hline
	\endfirsthead
	\caption*{Продолжение таблицы \ref{basekeymap:layouttable}}\\
	\hline \centrow 1 & \centrow 2 & \centrow 3\\ \hline
	\finishhead
	1 & Кнопки & Каждая из кнопок нужна при настройке той или иной функции клавиатуры
\end{xltabular}
\renewcommand{\arraystretch}{1.0} % восстановление сетки

На рисунке \ref{eurojapankeymap:layoutimage} изображена страница европейской и японской раскладок клавиатур.

\begin{figure}[H]
	\centering
	\includegraphics[width=1\linewidth]{eurojapankeymaplayout}
	\caption{Страница европейской и японской раскладок клавиатур}
	\label{eurojapankeymap:layoutimage}
\end{figure}

В таблице \ref{eurojapankeymap:layouttable} представлено описание объектов данного интерфейса.

\renewcommand{\arraystretch}{0.8} % уменьшение расстояний до сетки таблицы
\begin{xltabular}{\textwidth}{|p{2.25cm}|p{2.25cm}|X|}
	\caption{Описание объектов интерфейса страницы европейской и японской раскладок клавиатур\label{eurojapankeymap:layouttable}}\\
	\hline \centrow \setlength{\baselineskip}{0.7\baselineskip} Номер объекта & \centrow \setlength{\baselineskip}{0.7\baselineskip} Тип объекта & \centrow \setlength{\baselineskip}{0.7\baselineskip} Действие/описание объекта \\
	\hline \centrow 1 & \centrow 2 & \centrow 3\\ \hline
	\endfirsthead
	\caption*{Продолжение таблицы \ref{eurojapankeymap:layouttable}}\\
	\hline \centrow 1 & \centrow 2 & \centrow 3\\ \hline
	\finishhead
	1 & Кнопки & Каждая из кнопок нужна при настройке той или иной функции клавиатуры
\end{xltabular}
\renewcommand{\arraystretch}{1.0} % восстановление сетки

На рисунке \ref{printlayer:layoutimage} изображена страница слоёв для печати.

\begin{figure}[H]
	\centering
	\includegraphics[width=1\linewidth]{printlayerlayout}
	\caption{Страница слоёв для печати}
	\label{printlayer:layoutimage}
\end{figure}

В таблице \ref{printlayer:layouttable} представлено описание объектов данного интерфейса.

\renewcommand{\arraystretch}{0.8} % уменьшение расстояний до сетки таблицы
\begin{xltabular}{\textwidth}{|p{2.25cm}|p{2.25cm}|X|}
	\caption{Описание объектов интерфейса страницы слоёв для печати\label{printlayer:layouttable}}\\
	\hline \centrow \setlength{\baselineskip}{0.7\baselineskip} Номер объекта & \centrow \setlength{\baselineskip}{0.7\baselineskip} Тип объекта & \centrow \setlength{\baselineskip}{0.7\baselineskip} Действие/описание объекта \\
	\hline \centrow 1 & \centrow 2 & \centrow 3\\ \hline
	\endfirsthead
	\caption*{Продолжение таблицы \ref{printlayer:layouttable}}\\
	\hline \centrow 1 & \centrow 2 & \centrow 3\\ \hline
	\finishhead
	1 & Кнопки & Каждая из кнопок нужна при настройке той или иной функции клавиатуры
\end{xltabular}
\renewcommand{\arraystretch}{1.0} % восстановление сетки

На рисунке \ref{qmkkeys:layoutimage} изображена страница клавиш QMK.

\begin{figure}[H]
	\centering
	\includegraphics[width=1\linewidth]{qmkkeyslayout}
	\caption{Страница клавиш QMK}
	\label{qmkkeys:layoutimage}
\end{figure}

В таблице \ref{qmkkeys:layouttable} представлено описание объектов данного интерфейса.

\renewcommand{\arraystretch}{0.8} % уменьшение расстояний до сетки таблицы
\begin{xltabular}{\textwidth}{|p{2.25cm}|p{2.25cm}|X|}
	\caption{Описание объектов интерфейса страницы клавиш QMK\label{qmkkeys:layouttable}}\\
	\hline \centrow \setlength{\baselineskip}{0.7\baselineskip} Номер объекта & \centrow \setlength{\baselineskip}{0.7\baselineskip} Тип объекта & \centrow \setlength{\baselineskip}{0.7\baselineskip} Действие/описание объекта \\
	\hline \centrow 1 & \centrow 2 & \centrow 3\\ \hline
	\endfirsthead
	\caption*{Продолжение таблицы \ref{qmkkeys:layouttable}}\\
	\hline \centrow 1 & \centrow 2 & \centrow 3\\ \hline
	\finishhead
	1 & Кнопки & Каждая из кнопок нужна при настройке той или иной QMK функции клавиатуры
\end{xltabular}
\renewcommand{\arraystretch}{1.0} % восстановление сетки

На рисунке \ref{lightkeys:layoutimage} изображена страница клавиш подсветки.

\begin{figure}[H]
	\centering
	\includegraphics[width=1\linewidth]{lightkeyslayout}
	\caption{Страница клавиш подсветки}
	\label{lightkeys:layoutimage}
\end{figure}

В таблице \ref{lightkeys:layouttable} представлено описание объектов данного интерфейса.

\renewcommand{\arraystretch}{0.8} % уменьшение расстояний до сетки таблицы
\begin{xltabular}{\textwidth}{|p{2.25cm}|p{2.25cm}|X|}
	\caption{Описание объектов интерфейса страницы клавиш подсветки\label{lightkeys:layouttable}}\\
	\hline \centrow \setlength{\baselineskip}{0.7\baselineskip} Номер объекта & \centrow \setlength{\baselineskip}{0.7\baselineskip} Тип объекта & \centrow \setlength{\baselineskip}{0.7\baselineskip} Действие/описание объекта \\
	\hline \centrow 1 & \centrow 2 & \centrow 3\\ \hline
	\endfirsthead
	\caption*{Продолжение таблицы \ref{lightkeys:layouttable}}\\
	\hline \centrow 1 & \centrow 2 & \centrow 3\\ \hline
	\finishhead
	1 & Кнопки & Каждая из кнопок нужна при настройке подсветки клавиатуры с клавиатуры
\end{xltabular}
\renewcommand{\arraystretch}{1.0} % восстановление сетки

На рисунке \ref{mediamousekeys:layoutimage} изображена страница клавиш приложений, медиа и мыши.

\begin{figure}[H]
	\centering
	\includegraphics[width=1\linewidth]{mediamousekeyslayout}
	\caption{Страница клавиш приложений, медиа и мыши}
	\label{mediamousekeys:layoutimage}
\end{figure}

В таблице \ref{mediamousekeys:layouttable} представлено описание объектов данного интерфейса.

\renewcommand{\arraystretch}{0.8} % уменьшение расстояний до сетки таблицы
\begin{xltabular}{\textwidth}{|p{2.25cm}|p{2.25cm}|X|}
	\caption{Описание объектов интерфейса страницы клавиш приложений, медиа и мыши\label{mediamousekeys:layouttable}}\\
	\hline \centrow \setlength{\baselineskip}{0.7\baselineskip} Номер объекта & \centrow \setlength{\baselineskip}{0.7\baselineskip} Тип объекта & \centrow \setlength{\baselineskip}{0.7\baselineskip} Действие/описание объекта \\
	\hline \centrow 1 & \centrow 2 & \centrow 3\\ \hline
	\endfirsthead
	\caption*{Продолжение таблицы \ref{mediamousekeys:layouttable}}\\
	\hline \centrow 1 & \centrow 2 & \centrow 3\\ \hline
	\finishhead
	1 & Кнопки & Каждая из кнопок нужна при настройке как обычной клавиатуры, так и клавиатуры с трекболом 
\end{xltabular}
\renewcommand{\arraystretch}{1.0} % восстановление сетки

На рисунке \ref{tapdancekeys:layoutimage} изображена страница клавиш "<Tap Dance">.

\begin{figure}[H]
	\centering
	\includegraphics[width=1\linewidth]{tapdancekeyslayout}
	\caption{Страница клавиш "<Tap Dance">}
	\label{tapdancekeys:layoutimage}
\end{figure}

В таблице \ref{tapdancekeys:layouttable} представлено описание объектов данного интерфейса.

\renewcommand{\arraystretch}{0.8} % уменьшение расстояний до сетки таблицы
\begin{xltabular}{\textwidth}{|p{2.25cm}|p{2.25cm}|X|}
	\caption{Описание объектов интерфейса страницы клавиш "<Tap Dance">\label{tapdancekeys:layouttable}}\\
	\hline \centrow \setlength{\baselineskip}{0.7\baselineskip} Номер объекта & \centrow \setlength{\baselineskip}{0.7\baselineskip} Тип объекта & \centrow \setlength{\baselineskip}{0.7\baselineskip} Действие/описание объекта \\
	\hline \centrow 1 & \centrow 2 & \centrow 3\\ \hline
	\endfirsthead
	\caption*{Продолжение таблицы \ref{tapdancekeys:layouttable}}\\
	\hline \centrow 1 & \centrow 2 & \centrow 3\\ \hline
	\finishhead
	1 & Кнопки & Каждая из кнопок нужна при настройке "<Tap Dance"> и сохранения их профилей 
\end{xltabular}
\renewcommand{\arraystretch}{1.0} % восстановление сетки

На рисунке \ref{userkeys:layoutimage} изображена страница клавиш "<User">.

\begin{figure}[H]
	\centering
	\includegraphics[width=1\linewidth]{userkeyslayout}
	\caption{Страница клавиш "<User">}
	\label{userkeys:layoutimage}
\end{figure}

В таблице \ref{userkeys:layouttable} представлено описание объектов данного интерфейса.

\renewcommand{\arraystretch}{0.8} % уменьшение расстояний до сетки таблицы
\begin{xltabular}{\textwidth}{|p{2.25cm}|p{2.25cm}|X|}
	\caption{Описание объектов интерфейса страницы клавиш "<User">\label{userkeys:layouttable}}\\
	\hline \centrow \setlength{\baselineskip}{0.7\baselineskip} Номер объекта & \centrow \setlength{\baselineskip}{0.7\baselineskip} Тип объекта & \centrow \setlength{\baselineskip}{0.7\baselineskip} Действие/описание объекта \\
	\hline \centrow 1 & \centrow 2 & \centrow 3\\ \hline
	\endfirsthead
	\caption*{Продолжение таблицы \ref{userkeys:layouttable}}\\
	\hline \centrow 1 & \centrow 2 & \centrow 3\\ \hline
	\finishhead
	1 & Кнопки & Каждая из кнопок нужна при пользовательских настройках, которые индивидуальны для физиологии человека 
\end{xltabular}
\renewcommand{\arraystretch}{1.0} % восстановление сетки

На рисунке \ref{macrokeys:layoutimage} изображена страница клавиш "<Macro">.

\begin{figure}[H]
	\centering
	\includegraphics[width=1\linewidth]{macrokeyslayout}
	\caption{Страница клавиш "<Macro">}
	\label{macrokeys:layoutimage}
\end{figure}

В таблице \ref{macrokeys:layouttable} представлено описание объектов данного интерфейса.

\renewcommand{\arraystretch}{0.8} % уменьшение расстояний до сетки таблицы
\begin{xltabular}{\textwidth}{|p{2.25cm}|p{2.25cm}|X|}
	\caption{Описание объектов интерфейса страницы клавиш "<Macro">\label{macrokeys:layouttable}}\\
	\hline \centrow \setlength{\baselineskip}{0.7\baselineskip} Номер объекта & \centrow \setlength{\baselineskip}{0.7\baselineskip} Тип объекта & \centrow \setlength{\baselineskip}{0.7\baselineskip} Действие/описание объекта \\
	\hline \centrow 1 & \centrow 2 & \centrow 3\\ \hline
	\endfirsthead
	\caption*{Продолжение таблицы \ref{macrokeys:layouttable}}\\
	\hline \centrow 1 & \centrow 2 & \centrow 3\\ \hline
	\finishhead
	1 & Кнопки & Каждая из кнопок нужна при пользовательских настройках, которые индивидуальны для физиологии человека 
\end{xltabular}
\renewcommand{\arraystretch}{1.0} % восстановление сетки

